\section*{全册大纲}
\addcontentsline{toc}{section}{全册大纲}

本册给零基础用户补上“听得懂、写得出、能判断”的乐理底子，并尽量把每个概念落回夜半的钢琴卷帘、编曲区和混音台。它不是音乐学院的压缩教材，而是一条可自学的路线：先认音高，再数节奏，再听音程，最后能写出一段有功能的八小节。

\subsection*{读者与前置}

\begin{itemize}
  \item 不需要识谱，不要求会乐器；只需要能打开夜半、会点钢琴卷帘的绘制工具。
  \item 每章开头给出“学完能做什么”，结尾给出“在夜半里怎么做”与一组练习题；练习答案在章末。
  \item 术语第一次出现时用中文（英文），完整对照见附录 A；正文不堆同义词。
\end{itemize}

\subsection*{十二章路线}

\begin{enumerate}[label=\textbf{第 \arabic* 章}]
  \item \textbf{声音与音高}\\
  振动与频率、音名、八度、科学音高记法、MIDI 编号、十二平均律、中央 C 与 A4=440 Hz；把钢琴键盘、夜半音高行和频率公式对上。
  \item \textbf{节奏与节拍}\\
  脉冲与拍、拍号（2/4、3/4、4/4、6/8）、音符时值与休止、附点与连音、切分与反拍、速度与 BPM；在夜半的时间轴上数格子，用节拍器与走带对齐。
  \item \textbf{音程}\\
  半音与全音、大/小/增/减音程、单音程与复音程、协和与不协和、音程的听辨与唱名；在钢琴卷帘上量两个音符的行距。
  \item \textbf{音阶与调式}\\
  大调与小调、调号与升降记号、五声音阶、常用调式（多利亚、利底亚、混合利底亚等）、关系大小调；用夜半的音阶/和弦辅助工具检查音级。
  \item \textbf{和弦}\\
  三和弦、七和弦、转位、和弦标记（根音＋性质＋转位）、常用进行（I–V–vi–IV、ii–V–I）；在卷帘里叠音符，在混音台听低频与声部排列。
  \item \textbf{和声进行与功能}\\
  主/属/下属功能、终止式（正格、变格、半、阻碍）、借用和弦、副属和弦、转调入门；把“为什么这个和弦接那个”讲成可操作的规则。
  \item \textbf{旋律写作}\\
  动机、乐句、问答句、重复与模进、旋律与和弦音/张力音的关系、节奏动机；在夜半写四小节问答句并做两次变奏。
  \item \textbf{曲式入门}\\
  段落与层次、重复与变化、常见流行/民谣结构（AABA、 verse–chorus、intro–verse–chorus–bridge–outro）、段落过渡；用编曲区的片段复制与命名整理结构。
  \item \textbf{读谱与记谱}\\
  五线谱基础、谱号、音符与休止记号、表情与力度记号、简谱与 TAB 的存在意义；把五线谱位置与钢琴卷帘行号互相对照，能读简单谱例。
  \item \textbf{蓝调、五声与调式写作}\\
  五声的音集用法、蓝调音阶与蓝音、十二小节蓝调骨架与替换、调式特征音与调式循环；把“少两个音”写成风格，而不是简化版大调。
  \item \textbf{和声进阶：七和弦、声部进行与再配和}\\
  七和弦五种性质、解决七音与导音、平行五/八度、转位与低音线、drop-2、功能替换、持续低音与模进式再配和；从“能配和弦”走到“能改配和弦”。
  \item \textbf{节奏、乐句与曲式进阶}\\
  6/8 与 3/4、hemiola、三对二、切分与留白、动机四种发展、句眼与终止、段落功能与 16 小节结构；把节奏当结构工具，而不是只当鼓点。
\end{enumerate}

\subsection*{每章固定结构}

\begin{description}[leftmargin=2em,style=nextline]
  \item[概念] 用生活经验引入，先给直觉，再给定义与公式。
  \item[谱例/图示] TikZ 绘制的键盘、卷帘网格、五线谱或简谱示意；不依赖外部图片。
  \item[在夜半里怎么做] 只写活尖端代码与应用内手册中真实存在的入口；未提供的功能明确标“尚未提供”。
  \item[练习] 4–6 题，含答案或提示；至少一题必须在夜半里动手完成。
  \item[帮助主题映射] 本章小节 → 应用内帮助锚点（英文小写连字符）。
\end{description}

\subsection*{建议学习节奏}

\begin{itemize}
  \item 零基础：每周 1 章，先读完概念，再把练习在夜半里做一遍；第 1–5 章不要跳。
  \item 有乐器基础但不会数字音频：第 1、2、9 章快速过，重点做第 5–7 章的写作练习。
  \item 只想补“怎么用夜半写歌”：按第 1 → 2 → 5 → 7 → 8 章走，第 3、4、6 章当工具书查。
\end{itemize}

\subsection*{与其它分册的边界}

\begin{itemize}
  \item 本册讲“音乐为什么这样写”，不讲夜半每个菜单的完整操作；界面操作以《夜半用户手册》为准。
  \item 乐器音色、合成器参数与效果器录音链，放到《乐器与合成器手册》；本册只在配器与混音处点到为止。
  \item 编曲配器的工程化流程（轨道分层、总线、发送、母带），放到《编曲与配器手册》；本册第 6–12 章只给写作层面的结构。
\end{itemize}
